% !PRINTED 436
\weberpage{436}
If we now seek the condition that $\Theta^4 = 1$,
we have to set $\Theta^2 = \Theta^{-2}$. If, then, we form $\Theta^2$ and $\Theta^{-2}$
according to (10), one may first try to bring about the agreement
by setting one of the three numbers $\alpha, \beta, \gamma$, say $\gamma = 0$;
but then by (6) and (7) a second, say
$\beta = 0$, must also vanish, and the third must be $\alpha = -1$. But then $\Theta^2$ does not
agree with $\Theta^{-2}$. Hence all three quantities $\alpha, \beta, \gamma$
are different from zero. If we now set the second and third
terms of the first row in the substitutions
$\Theta^2$ and $\Theta^{-2}$ derived from (10) equal to one another, we obtain, when the factors
$\gamma$ and $\alpha$ are discarded:
\[
\alpha(\varepsilon - \varepsilon^{-2}) = \beta(\varepsilon^{-1} - 1), \quad \gamma(\varepsilon^4 - 1) = \beta(\varepsilon^3 - \varepsilon)
\]
or
\[
\alpha(\varepsilon^2 - \varepsilon^{-2}) = \beta(\varepsilon^4 - \varepsilon^{-4}), \quad \gamma(\varepsilon^2 - \varepsilon^{-2}) = \beta(\varepsilon - \varepsilon^{-1}),
\]
and from this, if $h$ denotes an undetermined factor:
\weq{11}{\begin{aligned} \alpha &= h(\varepsilon^4 - \varepsilon^{-4}) \\ \beta &= h(\varepsilon^2 - \varepsilon^{-2}) \\ \gamma &= h(\varepsilon - \varepsilon^{-1}). \end{aligned}}

The factor $h$ is determined according to (7) from the equation
\weq{12}{-1 = h(\varepsilon + \varepsilon^2 + \varepsilon^4 - \varepsilon^{-1} - \varepsilon^{-2} - \varepsilon^{-4})}
and by Bd. I, \S.\,171 we obtain
\weq{13}{h = \frac{i}{\sqrt{7}} = \frac{-1}{\sqrt{-7}}.}

The sign of $\sqrt{7}$ depends on the choice of $\varepsilon$ and
is positive if, e.\,g.,
\[
\varepsilon = e^{\frac{2\pi i}{7}}
\]
is taken. One then obtains for $\alpha, \beta, \gamma$
\[
\alpha = \frac{-2\sin\dfrac{8\pi}{7}}{\sqrt{7}}, \quad \beta = \frac{-2\sin\dfrac{4\pi}{7}}{\sqrt{7}}, \quad \gamma = \frac{-2\sin\dfrac{2\pi}{7}}{\sqrt{7}}.
\]

From (11) and (12) there follow further the formulae
\weq{14}{\begin{aligned}
\alpha\varepsilon + \beta\varepsilon^4 + \gamma\varepsilon^2 &= 1 \\
\alpha\varepsilon^{-1} + \beta\varepsilon^{-4} + \gamma\varepsilon^{-2} &= 1, \\
\alpha + \beta\varepsilon^{-1} + \gamma\varepsilon^2 &= 0 \\
\alpha + \beta\varepsilon + \gamma\varepsilon^{-2} &= 0; \\
\alpha\beta\gamma &= \tfrac{1}{7},
\end{aligned}}
%CONT>>
